\section{Monotonicity of Optimal Packing under the Instance Partial Order}
\label{app:partial-order-proof}

\begin{lemma}[Monotonicity of Optimal Packing]
\label{lem:monotonicity}
Let $I=\{a_1,\ldots,a_n\}$ and
$J=\{b_1,\ldots,b_m\}$ be two bin packing
instances. If $I\preceq J$, then
$\OPT(I)\leq\OPT(J).$
\end{lemma}


\begin{proof}
By the definition of $I\preceq J$, there exists
an injective mapping
$f:\{1,\ldots,n\}\to\{1,\ldots,m\}$
such that $a_i\leq b_{f(i)}$ for every
$i\in\{1,\ldots,n\}$.

Consider an optimal packing of $J$ using
$\OPT(J)$ bins. For each item $a_i\in I$,
replace the corresponding item $b_{f(i)}$
with $a_i$, and remove all unmatched items
from $J$.

Since $a_i\leq b_{f(i)}$, the total size
of items in each bin cannot increase.
Hence, the resulting packing of $I$ is
feasible and uses at most $\OPT(J)$ bins.

Therefore,
\[
    \OPT(I)\leq\OPT(J).
\]
\end{proof}

